\subsection{Letting a language model decide}
\label{sec:llmarms}

Everything so far concerns policies whose decisions are made by arithmetic. But the
systems this paper is about delegate the decision to a language model, and a critique of
that architecture that never runs it would not be worth much. This section runs it.

Three arms share the counterfactual tensor with every policy above, so the outcome of
each action is identical and the only thing that varies is who chooses.
\texttt{llm\_single} sees the diagnostic evidence once and picks an action, the
Self-Refine pattern. \texttt{llm\_critique\_gate} sees the same evidence but first judges
\emph{whether} the forecast needs revising at all: the reflection step whose value is
this paper's question, and the mechanism \citet{huang2024selfcorrect} and
\citet{kamoi2024selfcorrect} found unreliable in reasoning tasks.
\texttt{rgc\_llm\_router} hands the model only the \emph{which}, keeping our quantitative
gate for the \emph{whether}, which separates the value of gating from the value of
routing. All three receive exactly the evidence our own gate receives: the comparison is
about how the decision is made, not about who saw more.

\paragraph{How this section relates to the previous one.}
Section~\ref{sec:decomposition} found no realizable per-series headroom: choosing
repairs on validation does not beat accepting. That sets the bar for a language model
in the decision loop. It cannot win by choosing repairs cleverly, because there is
little to win; it can only avoid losing by declining to act. The question this section
asks is therefore whether a model restrains itself, and the answer is that it does not.

\paragraph{Three model families, because one would not settle it.}
A single model failing to beat never-correcting invites the obvious reply that we picked
a weak one. The arms therefore run on unrelated families: \texttt{gemini-3.5-flash-lite},
a small hosted proprietary model, on one instance from each of the \nseries{} series;
\texttt{qwen3.8-27b}, an open-weight model of different lineage served by a different
provider, on 140 series; and \texttt{gpt-oss-20b}, an open-weight \emph{reasoning} model
run at medium effort, on \reasonSeries{} series. Free-tier quotas set the coverage and they meter differently ---
one provider caps requests per day, the other tokens per day --- so the series count is
given per row in Table~\ref{tab:llmarms}, and the smaller family covers the two arms that
carry the argument rather than all three. Every non-LLM arm is re-scored within its own
family's subsample, so rows are comparable down a block but not across blocks.
Temperature is 0 throughout and every response is cached, so the experiment replays
without an API key.

The first two are not reasoning models and receive no prompt engineering: the evidence
block is a fixed template and each decision is a single zero-shot turn. The third
deliberates explicitly, and a fourth condition described below shows a model worked
examples drawn from the validation split. What the families establish together is
therefore not a claim about language models in general, but that the result is not an
artefact of one model's idiosyncrasies, of an absence of deliberation, or of the
zero-shot setting, which are the three objections a single-family null cannot answer.

\begin{table}[htbp]
\centering
\caption{LLM correction arms against the same counterfactual outcomes, for two model
families. Non-LLM arms are re-scored within each family's own subsample, so rows are
comparable within a family block but not across blocks; the covered series count is given
per row. Effects are series-level differences from never correcting with 95\% bootstrap
intervals; $\dagger$ marks an interval excluding zero. Negative favours the policy.
``LLM calls'' is the live provider call count, since a correction stage that does not
help is still not free.}
\label{tab:llmarms}
\small
\fittable{\begin{tabular}{lllrrrrrr}
\toprule
Metric & Model & Policy & Series & Series mean & $\Delta$ & 95\% CI & Win rate & LLM calls \\
\midrule
Point (MASE) & Gemini 3.5 Flash-Lite & never correct & 578 & 1.2220 & 0.00\% & [0.00, 0.00] & 50.0\% & 0 \\
Point (MASE) & Gemini 3.5 Flash-Lite & always ensemble & 578 & 1.1520 & -5.73\%$^\dagger$ & [-9.54, -1.99] & 53.9\% & 0 \\
Point (MASE) & Gemini 3.5 Flash-Lite & rgc calibrated & 578 & 1.2196 & -0.19\% & [-0.72, 0.26] & 50.3\% & 0 \\
Point (MASE) & Gemini 3.5 Flash-Lite & llm single & 578 & 1.3243 & 8.37\%$^\dagger$ & [3.29, 13.94] & 42.2\% & 0 \\
Point (MASE) & Gemini 3.5 Flash-Lite & llm critique gate & 578 & 1.3136 & 7.50\%$^\dagger$ & [2.73, 12.89] & 42.8\% & 288 \\
Point (MASE) & Gemini 3.5 Flash-Lite & rgc llm router & 578 & 1.2814 & 4.87\%$^\dagger$ & [0.95, 10.06] & 47.5\% & 258 \\
Point (MASE) & Gemini 3.5 Flash-Lite & always correct & 578 & 1.2920 & 5.73\% & [-0.48, 14.43] & 47.1\% & 0 \\
Probabilistic (CRPS) & Gemini 3.5 Flash-Lite & never correct & 578 & 1.0629 & 0.00\% & [0.00, 0.00] & 50.0\% & 0 \\
Probabilistic (CRPS) & Gemini 3.5 Flash-Lite & always ensemble & 578 & 0.9503 & -10.59\%$^\dagger$ & [-15.14, -6.70] & 53.8\% & 0 \\
Probabilistic (CRPS) & Gemini 3.5 Flash-Lite & rgc calibrated & 578 & 1.0064 & -5.31\%$^\dagger$ & [-9.36, -1.91] & 53.5\% & 0 \\
Probabilistic (CRPS) & Gemini 3.5 Flash-Lite & llm single & 578 & 1.1530 & 8.48\%$^\dagger$ & [3.00, 14.55] & 43.9\% & 0 \\
Probabilistic (CRPS) & Gemini 3.5 Flash-Lite & llm critique gate & 578 & 1.1493 & 8.13\%$^\dagger$ & [3.17, 13.92] & 44.3\% & 288 \\
Probabilistic (CRPS) & Gemini 3.5 Flash-Lite & rgc llm router & 578 & 1.1015 & 3.63\%$^\dagger$ & [0.48, 8.04] & 48.7\% & 258 \\
Probabilistic (CRPS) & Gemini 3.5 Flash-Lite & always correct & 578 & 1.0979 & 3.29\% & [-4.92, 13.64] & 51.6\% & 0 \\
Point (MASE) & Qwen3.8-27B (Groq) & never correct & 140 & 0.8936 & 0.00\% & [0.00, 0.00] & 50.0\% & 0 \\
Point (MASE) & Qwen3.8-27B (Groq) & always ensemble & 140 & 0.8892 & -0.50\% & [-5.17, 4.69] & 51.4\% & 0 \\
Point (MASE) & Qwen3.8-27B (Groq) & rgc calibrated & 140 & 0.8850 & -0.96\% & [-2.85, 0.00] & 50.7\% & 0 \\
Point (MASE) & Qwen3.8-27B (Groq) & llm single & 140 & 0.9400 & 5.19\%$^\dagger$ & [0.86, 10.19] & 46.8\% & 0 \\
Point (MASE) & Qwen3.8-27B (Groq) & llm critique gate & 140 & 0.9458 & 5.85\%$^\dagger$ & [1.45, 10.84] & 45.7\% & 11 \\
Probabilistic (CRPS) & Qwen3.8-27B (Groq) & never correct & 140 & 0.7513 & 0.00\% & [0.00, 0.00] & 50.0\% & 0 \\
Probabilistic (CRPS) & Qwen3.8-27B (Groq) & always ensemble & 140 & 0.7168 & -4.59\% & [-10.64, 0.96] & 52.1\% & 0 \\
Probabilistic (CRPS) & Qwen3.8-27B (Groq) & rgc calibrated & 140 & 0.7097 & -5.54\%$^\dagger$ & [-10.73, -0.69] & 57.1\% & 0 \\
Probabilistic (CRPS) & Qwen3.8-27B (Groq) & llm single & 140 & 0.7941 & 5.70\% & [-0.48, 12.33] & 48.6\% & 0 \\
Probabilistic (CRPS) & Qwen3.8-27B (Groq) & llm critique gate & 140 & 0.7892 & 5.04\% & [-1.22, 11.81] & 50.7\% & 11 \\
Point (MASE) & GPT-OSS-20B reasoning (Groq) & never correct & 60 & 1.0447 & 0.00\% & [0.00, 0.00] & 50.0\% & 0 \\
Point (MASE) & GPT-OSS-20B reasoning (Groq) & always ensemble & 60 & 1.0974 & 5.04\% & [-9.06, 24.71] & 48.3\% & 0 \\
Point (MASE) & GPT-OSS-20B reasoning (Groq) & rgc calibrated & 60 & 1.0454 & 0.06\% & [-0.16, 0.34] & 50.0\% & 0 \\
Point (MASE) & GPT-OSS-20B reasoning (Groq) & llm single & 60 & 1.1024 & 5.52\% & [-2.92, 14.61] & 50.0\% & 60 \\
Point (MASE) & GPT-OSS-20B reasoning (Groq) & llm critique gate & 60 & 1.1034 & 5.62\% & [-3.77, 15.32] & 47.5\% & 60 \\
Probabilistic (CRPS) & GPT-OSS-20B reasoning (Groq) & never correct & 60 & 0.9179 & 0.00\% & [0.00, 0.00] & 50.0\% & 0 \\
Probabilistic (CRPS) & GPT-OSS-20B reasoning (Groq) & always ensemble & 60 & 0.8967 & -2.31\% & [-16.91, 13.54] & 53.3\% & 0 \\
Probabilistic (CRPS) & GPT-OSS-20B reasoning (Groq) & rgc calibrated & 60 & 0.9341 & 1.76\% & [-3.18, 7.54] & 51.7\% & 0 \\
Probabilistic (CRPS) & GPT-OSS-20B reasoning (Groq) & llm single & 60 & 0.9776 & 6.49\% & [-2.30, 15.67] & 42.5\% & 60 \\
Probabilistic (CRPS) & GPT-OSS-20B reasoning (Groq) & llm critique gate & 60 & 0.9805 & 6.81\% & [-1.72, 15.52] & 43.3\% & 60 \\
Point (MASE) & Gemini 3.5 Flash-Lite, 8-shot & never correct & 60 & 1.0447 & 0.00\% & [0.00, 0.00] & 50.0\% & 0 \\
Point (MASE) & Gemini 3.5 Flash-Lite, 8-shot & always ensemble & 60 & 1.0974 & 5.04\% & [-9.06, 24.71] & 48.3\% & 0 \\
Point (MASE) & Gemini 3.5 Flash-Lite, 8-shot & rgc calibrated & 60 & 1.0454 & 0.06\% & [-0.16, 0.34] & 50.0\% & 0 \\
Point (MASE) & Gemini 3.5 Flash-Lite, 8-shot & llm single & 60 & 1.1142 & 6.65\% & [-1.25, 15.77] & 45.8\% & 0 \\
Point (MASE) & Gemini 3.5 Flash-Lite, 8-shot & llm critique gate & 60 & 1.3462 & 28.85\% & [-0.49, 80.30] & 46.7\% & 0 \\
Probabilistic (CRPS) & Gemini 3.5 Flash-Lite, 8-shot & never correct & 60 & 0.9179 & 0.00\% & [0.00, 0.00] & 50.0\% & 0 \\
Probabilistic (CRPS) & Gemini 3.5 Flash-Lite, 8-shot & always ensemble & 60 & 0.8967 & -2.31\% & [-16.91, 13.54] & 53.3\% & 0 \\
Probabilistic (CRPS) & Gemini 3.5 Flash-Lite, 8-shot & rgc calibrated & 60 & 0.9341 & 1.76\% & [-3.18, 7.54] & 51.7\% & 0 \\
Probabilistic (CRPS) & Gemini 3.5 Flash-Lite, 8-shot & llm single & 60 & 0.9944 & 8.33\%$^\dagger$ & [0.49, 16.66] & 40.0\% & 0 \\
Probabilistic (CRPS) & Gemini 3.5 Flash-Lite, 8-shot & llm critique gate & 60 & 1.1299 & 23.09\%$^\dagger$ & [1.96, 58.52] & 39.2\% & 0 \\
Point (MASE) & Gemini 3.5 Flash-Lite, 0-shot (matched) & never correct & 60 & 1.0447 & 0.00\% & [0.00, 0.00] & 50.0\% & 0 \\
Point (MASE) & Gemini 3.5 Flash-Lite, 0-shot (matched) & always ensemble & 60 & 1.0974 & 5.04\% & [-9.06, 24.71] & 48.3\% & 0 \\
Point (MASE) & Gemini 3.5 Flash-Lite, 0-shot (matched) & rgc calibrated & 60 & 1.0454 & 0.06\% & [-0.16, 0.34] & 50.0\% & 0 \\
Point (MASE) & Gemini 3.5 Flash-Lite, 0-shot (matched) & llm single & 60 & 1.1330 & 8.45\%$^\dagger$ & [3.00, 14.72] & 39.2\% & 50 \\
Point (MASE) & Gemini 3.5 Flash-Lite, 0-shot (matched) & llm critique gate & 60 & 1.1407 & 9.19\%$^\dagger$ & [1.69, 17.60] & 38.3\% & 50 \\
Probabilistic (CRPS) & Gemini 3.5 Flash-Lite, 0-shot (matched) & never correct & 60 & 0.9179 & 0.00\% & [0.00, 0.00] & 50.0\% & 0 \\
Probabilistic (CRPS) & Gemini 3.5 Flash-Lite, 0-shot (matched) & always ensemble & 60 & 0.8967 & -2.31\% & [-16.91, 13.54] & 53.3\% & 0 \\
Probabilistic (CRPS) & Gemini 3.5 Flash-Lite, 0-shot (matched) & rgc calibrated & 60 & 0.9341 & 1.76\% & [-3.18, 7.54] & 51.7\% & 0 \\
Probabilistic (CRPS) & Gemini 3.5 Flash-Lite, 0-shot (matched) & llm single & 60 & 0.9981 & 8.73\%$^\dagger$ & [2.75, 15.62] & 38.3\% & 50 \\
Probabilistic (CRPS) & Gemini 3.5 Flash-Lite, 0-shot (matched) & llm critique gate & 60 & 1.0195 & 11.06\%$^\dagger$ & [4.34, 18.85] & 34.2\% & 50 \\
\bottomrule
\end{tabular}
}
\end{table}

\paragraph{On the primary family, every LLM arm is worse than doing nothing.}
On the primary family all
three arms degrade both metrics with intervals that exclude zero: a single pass costs
\llmMaseLlmSinglePct\% MASE ([\llmMaseLlmSingleCIlo, \llmMaseLlmSingleCIhi]\%) and
\llmCrpsLlmSinglePct\% CRPS ([\llmCrpsLlmSingleCIlo, \llmCrpsLlmSingleCIhi]\%); the
self-critique gate costs \llmMaseLlmCritiqueGatePct\% and
\llmCrpsLlmCritiqueGatePct\%, and even the router, which acts on only a seventh of
instances because our quantitative gate holds it back, costs
\llmMaseRgcLlmRouterPct\% and \llmCrpsRgcLlmRouterPct\%. Each wins on fewer than half of
series. On the same instances the static ensemble improves CRPS by
\llmCrpsAlwaysEnsembleAbs\% ([\llmCrpsAlwaysEnsembleAbsCIlo,
\llmCrpsAlwaysEnsembleAbsCIhi]\%), so the subsample is not too small to resolve a real
effect; it resolves one, in the opposite direction from the arms that reason about each
instance.

\paragraph{A discrepancy between two designs.}
In Table~\ref{tab:llmarms} the static ensemble's MASE advantage clears significance,
whereas in Table~\ref{tab:headline} the same policy's MASE interval includes zero. The
tables are computed on different data: the headline analysis uses every test origin
(\ntestinstances{} instances), while these arms use one randomly chosen origin per series,
because that is what the API budget allowed. Both are series-level and both are correct
for what they cover; the subsample simply draws fewer of the extreme origins that widen
the full interval. \textbf{The headline analysis is the primary one}, and the claim we
stand behind is the more conservative of the two: that no policy reliably improves
point accuracy. We flag the difference because quoting the subsample's narrower interval
as though it settled the point would be exactly the kind of selection this paper
criticises elsewhere.

The decision-level view explains the aggregate. When the language model chooses to
correct, it makes the forecast \emph{worse} on roughly half the instances it touches and
better on under a third: it harms about $1.8\times$ as often as it helps. The router
is the worst offender at a $64\%$ harm rate, which is initially surprising until one
notices what it means: even after a calibrated gate has correctly identified an instance
worth correcting, delegating the choice of \emph{which} repair to a language model
destroys the opportunity. That is Section~\ref{sec:decomposition}'s restriction-and-first-choice
loss appearing again, with a language model in the role of the fixed canonical mapping.
A better-informed guess is still a guess; it is not a measurement.

\paragraph{The reflection step carries almost no information.}
We expected the comparison between \texttt{llm\_single} and \texttt{llm\_critique\_gate}
to be the sharpest in this section, because they differ in exactly one thing: whether the
model is asked to judge the need for correction before choosing. It is not sharp at all.
Asked plainly whether the forecast needs revising, the model answers yes on
\rateLlmCritiqueGate\% of instances, almost exactly the rate at which the single-pass arm
chooses to act anyway, and the two arms land within a percentage point of each other on
both metrics.

Matching rates would be weak evidence on its own; two arms can intervene equally often
while selecting very different instances, one of them well. So we test discrimination
directly, against the benchmark's ground-truth \texttt{should\_correct} label
(Table~\ref{tab:discrimination}). The critique's yes/no separates warranted from
unwarranted correction at AUROC \discrimAurocCritique{}, against a chance value of 0.500.
Its selection precision is \discrimPrecisionCritique\%, where the base rate is
\discrimBaseRate\% and a coin firing at the same rate achieves \discrimPrecisionRandom\%.
And \discrimJaccard\% of the instances it elects to correct are the same instances the
uncritiqued arm corrects.

The decisive comparison is with the arm that does no critique at all, which discriminates
at AUROC \discrimAurocSingle{}. \textbf{Adding an explicit reflection step moves
discrimination by \discrimAurocDelta{}; from near-chance to near-chance.} The critique
is not merely a weak gate; it is very nearly a re-description of what the model would have
done without being asked. Both families show the same pattern
(Table~\ref{tab:discrimination}), including one where the critique step \emph{raises} the
intervention rate instead of leaving it flat, and still adds essentially nothing to
discrimination.

\paragraph{The same judgement, made arithmetically.}
The natural question is whether this decision is simply hard. It is not. Our reliability
estimator --- gradient boosting over the same diagnostic features the model was shown in
prose, fitted on validation only --- scores \discrimAurocArithmetic{} against the
\emph{same} label on the \emph{same} instances (Table~\ref{tab:discrimination}). That is
\discrimArithmeticAdvantage{} AUROC above the self-critique. We deliberately do not
express that as a multiple of the critique's distance from chance: dividing by a
near-zero baseline is the first of the measurement pathologies this paper documents
(Section~\ref{sec:pathologies}). We also
report the contrast \emph{matched}, because the triage figure quoted in
Section~\ref{sec:triage} targets a different label, and setting the two against each
other directly would not be legitimate.

\paragraph{What that contrast does and does not license.}
A binary yes/no scored as an AUROC is balanced accuracy, and nothing more. The
estimator emits a continuous score, so scoring the two the same way credits it for
ranking that a yes/no cannot express, and some of the apparent gap is an artefact of
the output type rather than of the judgement behind it. The contrast is therefore reported at matched footing in two ways, both in
Table~\ref{tab:discrimination}.

The first gives the language model the same kind of credit. Its verbalised confidence
is continuous, so it can be scored as a ranking against the same label: it reaches
AUROC \discrimAurocLlmConfidence{} against the estimator's \discrimAurocArithmetic{},
a gap of \discrimContinuousGap{}. The second removes the ranking question entirely
by forcing both to the same operating point --- the estimator thresholded to flag
exactly the fraction of instances the critique gate flags --- and compares balanced
accuracy: \discrimBalaccArithmetic{} against \discrimBalaccCritique{}, a gap of
\discrimMatchedGap{}.

The two ask different questions and both are worth keeping. (a) asks whether the
model's stated risk can \emph{rank} instances at all, independent of where a threshold
would be placed. (b) asks how much better the decision is at the rate the model
actually fires, which is the quantity a deployed gate delivers. It is worth noting that
(a) is marginally \emph{larger} than the unmatched comparison it replaces
(\discrimContinuousGap{} against \discrimArithmeticAdvantage{}), so matching was not a
route to a smaller number; only (b), which removes the ranking question entirely,
narrows the gap.

At a matched operating point the gap
is \discrimMatchedGap{}, not the \discrimArithmeticAdvantage{} a naive reading of the
first two columns would give. The artefact was real and it was worth roughly a third
of the apparent difference. What survives it is still the finding: on the same
instances, against the same label, at the same rate of intervention, a gradient-boosted
model over a few dozen scalar diagnostics decides substantially better than a language
model reasoning over the same evidence in words.

What the contrast does not license is a claim that language models cannot do this. The
estimator is fitted on the validation split and the model is not, and that asymmetry is
addressed directly below. Nor does it license a claim about relative ceilings: it
compares one fitted estimator against one prompted model, at one budget, and a
different prompt or a different model could land elsewhere. It licenses exactly one
thing, which is the thing the deployment question turns on: that the cheap
arithmetic option available to any practitioner today is the better of the two.


\begin{table}[htbp]
\centering
\caption{Does verbal self-critique carry information about whether to correct? All
columns are computed on the same instances against the same ground-truth
\texttt{should\_correct} label. ``AUROC (critique)'' scores a binary yes/no and is
therefore balanced accuracy; ``AUROC (arithmetic)'' scores a continuous estimator and
is credited for ranking, so the two are compared at matched footing in the text rather
than by subtracting these columns. $\Delta$ AUROC is the within-model quantity the
claim rests on: the same model, the same instances, deciding with and without being
asked to reflect. ``Overlap'' is the Jaccard overlap of the two intervention sets.}
\label{tab:discrimination}
\small
\fittable{\begin{tabular}{lrrrrrrrrr}
\toprule
Model & $n$ & Base rate & AUROC (critique) & AUROC (no critique) & $\Delta$ AUROC & AUROC (arithmetic) & Precision & Random & Overlap \\
\midrule
Gemini 3.5 Flash-Lite & 578 & 59.7\% & 0.532 & 0.522 & 0.010 & 0.760 & 61.9\% & 59.9\% & 82\% \\
Qwen3.8-27B (Groq) & 140 & 52.9\% & 0.568 & 0.561 & 0.008 & 0.693 & 59.0\% & 51.4\% & 75\% \\
GPT-OSS-20B reasoning (Groq) & 60 & 60.0\% & 0.493 & 0.590 & -0.097 & 0.774 & 59.6\% & 60.0\% & 80\% \\
\bottomrule
\end{tabular}
}
\end{table}

\paragraph{Why it intervenes so often: it thinks its forecasts are far worse than
they are.}
The arms are asked, in the same breath as the decision, for a number between 0 and 1
expressing how reliable they judge the current forecast to be. That is a claim about
exactly the event the reliability estimator predicts, so the two can be scored against
each other on identical instances and an identical label. This is pre-registered
hypothesis H5, and it is \textbf{supported}: the estimator's calibration error is
\calEceDiff{} lower ([\calEceDiffCIlo, \calEceDiffCIhi], bootstrapped over series),
and the same ordering holds on the Brier score.

The size and direction of the miscalibration are what matter here. Large errors occur
on \calBaseRate\% (on the \calN-instance subsample the LLM arms run on) of these instances; the model's stated risk averages
\calLlmMeanRisk{}. It believes its forecasts fail roughly two and a half times as often
as they do. That is not a detail about calibration, it is the mechanism behind the
headline: a model that thinks most of its forecasts are unreliable will conclude that
most of them need fixing, which is precisely the disposition the intervention rates
show. The reflection step is not failing to detect a problem it should have seen: it is reporting a problem that is mostly not there.

Table~\ref{tab:calibration} also contains the reason we do not rest this on calibration
error alone. A constant predictor that simply emits the base rate attains an ECE of
essentially zero while ranking nothing at all, and its Brier score (\calConstBrier{})
is nonetheless better than the language model's (\calLlmBrier{}). By the proper
scoring rule, the agent's stated confidence is worse than a constant. It is not
worthless as a \emph{ranking}: its AURC of \calLlmAurc{} beats the constant's
\calConstAurc{}, so there is some signal in which forecasts it doubts most. But that
signal is buried under a bias large enough to make the resulting decisions wrong, and a
gate reads the level, not the ranking.

That a language model's stated confidence is poorly calibrated is not itself news
\citep{kadavath2022know,lin2022verbalized,xiong2024confidence}, and we do not present it
as such. Two things here are worth separating from that literature. The direction is
inverted: the usual finding is overconfidence, a model claiming more certainty than its
accuracy warrants, whereas this model is systematically \emph{pessimistic} about the
artefact it is judging. And the consequence is operational rather than presentational.
In the settings where verbalised confidence is normally studied it is a number shown to
a reader, who may discount it; here it is wired into a control loop, where a bias of
this size does not mislead anyone so much as it manufactures roughly one unnecessary
repair for every instance that needed one.

\begin{table}[htbp]
\centering
\caption{H5: the agent's stated confidence against a calibrated estimate of the same
event. Both predict the probability of a large error, on the same \calN{} instances
and against the same label; lower is better in every column. The constant predictor is
included deliberately: it attains a perfect calibration error while ranking nothing,
which is why ECE is never read here on its own.}
\label{tab:calibration}
\small
\fittable{\begin{tabular}{lrrrr}
\toprule
Predictor & ECE & Brier & AURC & Mean $p$ \\
\midrule
LLM verbalised confidence & 0.3254 & 0.3148 & 0.1593 & 0.544 \\
Reliability estimator (calibrated) & 0.0502 & 0.1401 & 0.0883 & 0.185 \\
Constant at base rate (0.22) & 0.0000 & 0.1714 & 0.2291 & 0.220 \\
\bottomrule
\end{tabular}
}
\end{table}

\paragraph{Deliberation, and demonstration, against criteria fixed in advance.}
Two reasonable objections are still open at this point. The models run so far answer in a
single zero-shot turn, so the failure could be an absence of deliberation instead of a
property of the mechanism. And the arithmetic comparator is fitted on the validation
split while the language model is shown nothing, so the contrast above could be measuring
supervision rather than capability. Both were tested on one shared denominator ---
\reasonSeries{} series, one origin each --- with the falsification criteria written down
before the runs completed (deviation~D9). The deliberation arm is \texttt{gpt-oss-20b} at
\texttt{reasoning\_effort} \texttt{= medium}, with a completion budget large enough for
the reasoning trace and the answer; it recorded no provider failures and no unparsed
replies. The demonstration arm prepends eight worked examples drawn from validation and
labelled by the counterfactual tensor --- the same supervision, out of the same data, that
the estimator receives --- and runs against a matched zero-shot control on the identical
series, so the comparison is within-model on the same instances.

Neither closes the gap. The reasoning model degrades MASE by \reasonMaseSinglePct\%
([\reasonMaseSingleCIlo, \reasonMaseSingleCIhi]\%) in one pass and
\reasonMaseCritiquePct\% ([\reasonMaseCritiqueCIlo, \reasonMaseCritiqueCIhi]\%) with
self-critique, intervening on \reasonRateSingle\% and \reasonRateCritique\% of instances
respectively: the highest rates of any family. Eight worked examples move the
intervention rate from \zeroshotRateSingle\% to \fewshotRateSingle\% and the effect on
MASE from \zeroshotMaseSinglePct\% ([\zeroshotMaseSingleCIlo,
\zeroshotMaseSingleCIhi]\%, excluding zero) to \fewshotMaseSinglePct\%
([\fewshotMaseSingleCIlo, \fewshotMaseSingleCIhi]\%, no longer excluding zero). That is a small improvement this denominator cannot resolve, not a repair: the
two intervals overlap over almost their whole length, and the demonstration condition does
not bring the arm to parity with never correcting. The critique gate is a different story, and it runs against us. With the same eight
examples its degradation grew rather than shrank: from \zeroshotMaseCritiquePct\% to
\fewshotMaseCritiquePct\% on MASE ([\fewshotMaseCritiqueCIlo, \fewshotMaseCritiqueCIhi]\%)
and from \zeroshotCrpsCritiquePct\% to \fewshotCrpsCritiquePct\% on CRPS
([\fewshotCrpsCritiqueCIlo, \fewshotCrpsCritiqueCIhi]\%, excluding zero). The MASE
interval is wide enough that a few series drive it, but the direction is unambiguous:
showing the gate worked examples made it more harmful, not less. We do not have an
explanation we can test. What it does rule out is the reading that the critique gate
failed only for want of supervision.

The obvious question is why the demonstration condition stopped at eight examples and
\fewshotSeries{} series. The answer is budget rather than design: the prefix multiplies
prompt tokens roughly sixfold, and free-tier quotas bind on tokens for one provider and
on requests for the other. A larger demonstration budget, or in-context examples
selected per instance rather than fixed, is untested here and is the version of this
arm we would run first with more compute. What we can say is that the cheapest form of
the intervention, the one a practitioner would reach for, does not restrain the
model, and that is a narrower claim than the arm was built to make.

The pre-registered falsifiers did not fire. The primary one asked whether self-critique
would raise this model's discrimination by more than $+0.05$ AUROC over the same model
deciding without reflection. It moved by \reasonAurocDelta{}: \emph{down}, from
\reasonAurocSingle{} to \reasonAurocCritique{}, on \reasonDiscrimN{} instances where the
arithmetic estimator scored \reasonAurocArithmetic{} against the identical label. This is
the strongest form of the within-model result in the paper. Asking the one model that
actually deliberates to deliberate about its own judgement did not improve that judgement;
it made it slightly worse.

\paragraph{Given the decision to correct, is the chosen repair any good?}
The arms conflate two decisions: whether to correct, and which repair to apply. The
tensor separates them, because it holds the outcome of every action on every instance.
On the instances where an arm intervened we compare the repair it chose against a
uniform draw from the same menu on the same instance, computed exactly by averaging the
menu's outcomes. That comparison involves no selection, so it is the one the claims rest
on (Table~\ref{tab:actionchoice}).

On the primary family the choice cannot be told apart from the coin. Single-pass
selection scores \actionChoiceSinglePct\% relative to a uniform draw
([\actionChoiceSingleCIlo, \actionChoiceSingleCIhi]\%) and the critique arm
\actionChoiceCritiquePct\% ([\actionChoiceCritiqueCIlo, \actionChoiceCritiqueCIhi]\%).
Having decided to act, the model's choice of repair carries no measurable advantage over
choosing at random among the options it was shown.

The table also shows the best action in each menu. That row is a maximum over test
outcomes, and Section~\ref{sec:decomposition} shows such maxima are almost entirely
winner's curse on this benchmark, so it is context rather than a target, and no claim
uses it. The router arm's \actionChoiceRouterPct\% rests on \actionChoiceRouterN{}
instances with an interval of [\actionChoiceRouterCIlo, \actionChoiceRouterCIhi]\%; it
is underpowered and excluded from claims.

\begin{table}[htbp]
\centering
\caption{Given the decision to correct, how good is the chosen repair? MASE on the
instances where each arm intervened. The uniform draw averages the menu's outcomes
exactly and is the comparison claims rest on. The best-in-menu row is a maximum over
test outcomes, inflated by selection (Section~\ref{sec:decomposition}), and is shown
for context only. Intervals are series-level bootstrap; the router arm is underpowered.}
\label{tab:actionchoice}
\small
\fittable{\begin{tabular}{lrrr}
\toprule
 & LLM single & LLM critique gate & RGC + LLM router \\
\midrule
Instances intervened on & 361 & 367 & 80 \\
Best in menu (test-selected ceiling) & 0.6487 & 0.6273 & 0.5067 \\
\textbf{Model's chosen action} & \textbf{1.1190} & \textbf{1.0906} & \textbf{1.4109} \\
Uniform draw from same menu & 1.0806 & 1.0678 & 1.0996 \\
Worst action in menu & 2.0608 & 2.0461 & 2.5194 \\
\midrule Chosen vs.\ uniform draw & 3.55\%~[-4.07, 12.11] & 2.13\%~[-5.17, 10.66] & 28.31\%~[-0.99, 65.84] \\
\bottomrule
\end{tabular}
}
\end{table}

\paragraph{The reasoning model may choose better, but it still acts too often.}
The reasoning family points the other way on this test. Its single-pass choice scores
\reasonChoiceSinglePct\% against a uniform draw ([\reasonChoiceSingleCIlo,
\reasonChoiceSingleCIhi]\%) and its critique arm \reasonChoiceCritiquePct\%
([\reasonChoiceCritiqueCIlo, \reasonChoiceCritiqueCIhi]\%). Both intervals just
include zero at \reasonSeries{} series, so this is a direction, not an effect. If it
holds, deliberation helps with which repair to apply and not with whether to apply one;
and because this arm still intervenes on \reasonRateSingle\% of instances, its net
effect on accuracy stays negative either way.

\paragraph{The pre-registered test of H7.}
H7 asked whether a quantitative gate decides better than verbal self-critique, scored by
oracle regret. It is the one hypothesis in the plan that could not be evaluated until
these arms ran. It is \textbf{\verdictHSeven}: the calibrated gate's mean regret is
\statHSeven{} lower than the LLM critique gate's. The comparison is worth stating
carefully, because the gate achieves this largely by declining to act; its regret is
almost identical to never correcting. The finding is therefore not that quantitative
gating is good at choosing corrections. It is that it is good at \emph{not} choosing
them, and that verbal self-critique is not.

\paragraph{What this does not establish.}
Three families, a reasoning model at real effort, and a demonstration condition close the
cheapest objections but not every one. We cannot rule out that a frontier reasoning model,
multi-turn deliberation with tool access during the decision, or a prompt engineered
against this benchmark would do better; nor that a larger member of the same reasoning
family --- the one originally planned, lost to a daily token cap --- would separate from
the 20B model we ran. The reasoning and demonstration conditions cover \reasonSeries{}
series, which is enough to replicate a direction and not enough to establish an effect on
its own. Nor did we test sensitivity to prompt wording, and the pessimism result in particular
may depend on how the confidence question is phrased. What the result does establish
is narrower: the mechanism
\emph{as commonly deployed}, on the evidence commonly available to it, is not merely
unhelpful but reliably harmful; its harm resolves at the primary family's sample size, while the two smaller families show the same direction without resolving it. Self-reflection, the intervention most often proposed as the fix, did not improve the model's judgement in any family we ran.

\paragraph{The cost side.}
These are also the expensive arms. The self-critique gate alone made
\llmCallsCritiqueGate{} live provider calls, and the three arms together consumed
\llmCallsTotal{} calls and \llmPromptTokensThousands{}k prompt tokens on the primary
family. The static ensemble that outperforms all of them makes zero
provider calls and takes no per-instance decision whatsoever. Any argument for an
LLM-in-the-loop correction stage has to be made against that baseline, and on this
benchmark it cannot be made on accuracy at any price.
